\DocumentMetadata{
  pdfversion=2.0,
  pdfstandard=ua-2,
  lang=en-US,
  tagging=on,
  tagging-setup={math/setup=mathml-SE,tabsorder=structure}
}
\documentclass[11pt,letterpaper]{article}
\usepackage[margin=0.68in,includefoot]{geometry}
\usepackage{newcomputermodern}
\usepackage{amsmath,amscd,mathtools}
\usepackage{tikz,adjustbox}
\providecommand{\square}{\mdlgwhtsquare}
\usepackage[hidelinks]{hyperref}
\hypersetup{
  pdftitle={Algebra First-Year Exam, January 2025, Part 2},
  pdfauthor={University of Florida Department of Mathematics},
  pdfsubject={Graduate mathematics examination},
  pdfdisplaydoctitle=true,
  bookmarksopen=true
}
\setcounter{secnumdepth}{0}
\setlength{\parindent}{0pt}
\setlength{\parskip}{0.28em}
\setlength{\emergencystretch}{3em}
\setlength{\leftmargini}{2.15em}
\setlength{\leftmarginii}{2.65em}
\setlength{\leftmarginiii}{3em}
\renewcommand{\labelenumi}{\textbf{\theenumi.}}
\pagestyle{empty}
\raggedbottom
\allowdisplaybreaks
\newenvironment{examproblems}[1]{%
  \begin{enumerate}%
  \setcounter{enumi}{#1}%
  \setlength{\topsep}{0.35em}%
  \setlength{\partopsep}{0pt}%
  \setlength{\itemsep}{0.48em}%
  \setlength{\parsep}{0.18em}%
}{\end{enumerate}}
\newenvironment{examsubparts}{%
  \begin{enumerate}%
  \setlength{\topsep}{0.18em}%
  \setlength{\partopsep}{0pt}%
  \setlength{\itemsep}{0.16em}%
  \setlength{\parsep}{0.1em}%
}{\end{enumerate}}
\newenvironment{examinstructions}{%
  \par\begingroup\itshape
}{%
  \par\endgroup\vspace{0.18em}
}
\makeatletter
\renewcommand\section{\@startsection{section}{1}{0pt}%
  {-1.25ex plus -.25ex minus -.1ex}%
  {0.55ex plus .1ex}%
  {\normalfont\large\bfseries}}
\renewcommand{\@maketitle}{%
  \newpage\null\vskip -1em%
  {\footnotesize\bfseries
    \href{https://www.ufl.edu/}{University of Florida}, \href{https://math.ufl.edu/}{Department of Mathematics}\par}%
  \vskip -0.5em%
  \tagstructbegin{tag=Title}%
    {\LARGE\bfseries\href{https://gma.math.ufl.edu/exams/algebra-first-year/}{Algebra First-Year Exam}, January 2025, Part 2\par}%
  \tagstructend%
  \vskip 0.25em%
  \tagmcbegin{artifact}%
    \hrule height 1pt%
  \tagmcend%
  \vskip 1em%
}
\makeatother
\title{Algebra First-Year Exam, January 2025, Part 2}
\author{}
\date{}
\begin{document}
\maketitle
\thispagestyle{empty}
\begin{examinstructions}
Answer four problems. (If you turn in more, the first four will be graded.) Within reason, you may quote theorems as long as you state them clearly.
\end{examinstructions}
\begin{examproblems}{0}
\item
Show that if the ring~\(R\) is a Euclidean domain, then~\(R\) is a principal ideal domain.
\item
Suppose that~\(F\) is a field and~\(G\) is a finite multiplicative subgroup of \(F\setminus\{0\}\). Prove that~\(G\) is cyclic.
\item
Suppose that~\(F\) is a field whose characteristic is not~\(2\). Assume that \(d_1,d_2\in F\) are not squares in~\(F\). Prove that \(F(\sqrt{d_1},\sqrt{d_2})\) is of dimension~\(4\) over~\(F\) if \(d_1d_2\) is not a square in~\(F\) and of dimension~\(2\) otherwise.
\item
Let \(GL_4(\mathbb{Z}/3\mathbb{Z})\) denote the group of all invertible~\(4\) by~\(4\) matrices with entries in \(\mathbb{Z}/3\mathbb{Z}\), the field of three elements. Use rational canonical forms to determine the number of conjugacy classes of elements of order~\(4\). Give the rational canonical form for each class.
\item
Let~\(R\) be an integral domain, and~\(M\) be an~\(R\)-module.
\begin{examsubparts}
\item[\textnormal{(a)}]
Define the rank of~\(M\). (Note that this definition is different from the definition of free module of some rank. In particular, it applies to modules that are not torsion free.)
\item[\textnormal{(b)}]
Prove that if~\(M\) is a free~\(R\)-module of rank~\(n\), (where~\(n\) is a non-negative integer), then the rank of any submodule of~\(M\) is at most~\(n\).
\end{examsubparts}
\item
Let~\(F\) be a field.
\begin{examsubparts}
\item[\textnormal{(a)}]
What does it mean to say that a field extension of~\(F\) is algebraic?
\item[\textnormal{(b)}]
Let \(F\subseteq K\subseteq L\) be three fields. Prove that if~\(K\) is algebraic over~\(F\) and~\(L\) is algebraic over~\(K\), then~\(L\) is algebraic over~\(F\).
\end{examsubparts}
\end{examproblems}
\end{document}
